% ============================================================================
% WORKBOOK — Section 3.7: Converting Rational Numbers to Decimals
% CCSS 7.NS.A.2.D
% Practice-focused reinforcement for the study guide topic
% ============================================================================
\section{Converting Rational Numbers to Decimals}
\topicTitle{Converting Rational Numbers to Decimals}

\begin{quickReview}{Terminating and Repeating Decimals}
Every rational number becomes a decimal that either \textbf{terminates} or \textbf{repeats}.

\begin{itemize}[leftmargin=*]
    \item \textbf{Terminating decimal:} The division comes out even. Example: $\frac{3}{8} = 0.375$.
    \item \textbf{Repeating decimal:} One or more digits repeat forever. Example: $\frac{1}{3} = 0.\overline{3}$.
    \item Use a \textbf{bar} over the repeating block: $\frac{2}{11} = 0.\overline{18}$ means $0.181818\ldots$
    \item \textbf{Quick test:} If the denominator (in lowest terms) has only factors of $2$ and $5$, the decimal terminates. Otherwise, it repeats.
\end{itemize}

\medskip
\textbf{Example:} $\frac{5}{6} = 0.8333\ldots = 0.8\overline{3}$ \quad (the $3$ repeats).
\end{quickReview}

% ── Warm-Up ──────────────────────────────────────────────────────────────────
\begin{practiceBox}[funGreen]{\faSun~Warm-Up}

\practiceHeader[funGreen]{Convert each fraction to a decimal.}

\begin{multicols}{2}
\prob $\dfrac{1}{2}$ \answerBlank[2cm]
\answerExplain{$0.5$}{Divide $1 \div 2 = 0.5$. The denominator is $2$ (a factor of $2$ only), so this terminates.}

\prob $\dfrac{3}{4}$ \answerBlank[2cm]
\answerExplain{$0.75$}{Divide $3 \div 4 = 0.75$. The denominator $4 = 2^2$ has only the factor $2$, so it terminates.}

\prob $\dfrac{1}{5}$ \answerBlank[2cm]
\answerExplain{$0.2$}{Divide $1 \div 5 = 0.2$. The denominator is $5$, so this terminates.}

\prob $\dfrac{1}{3}$ \answerBlank[2cm]
\answerExplain{$0.\overline{3}$}{Divide $1 \div 3 = 0.333\ldots$ The digit $3$ repeats forever: $0.\overline{3}$.}

\prob $\dfrac{2}{5}$ \answerBlank[2cm]
\answerExplain{$0.4$}{Divide $2 \div 5 = 0.4$. The denominator is $5$, so it terminates.}

\prob $\dfrac{7}{10}$ \answerBlank[2cm]
\answerExplain{$0.7$}{Divide $7 \div 10 = 0.7$. The denominator $10 = 2 \times 5$, so it terminates.}
\end{multicols}
\end{practiceBox}

% ── Practice A: Terminating Decimals ─────────────────────────────────────────
\begin{practiceBox}{Terminating Decimals}

\practiceHeader{Convert each fraction to a terminating decimal.}

\begin{multicols}{2}
\prob $\dfrac{7}{8}$ \answerBlank[2cm]
\answerExplain{$0.875$}{Divide $7 \div 8$: $8$ into $70$ is $8$ R$6$; $8$ into $60$ is $7$ R$4$; $8$ into $40$ is $5$ R$0$. So $\frac{7}{8} = 0.875$.}

\prob $\dfrac{3}{16}$ \answerBlank[2cm]
\answerExplain{$0.1875$}{Divide $3 \div 16 = 0.1875$. The denominator $16 = 2^4$ has only factor $2$, so it terminates.}

\prob $\dfrac{9}{20}$ \answerBlank[2cm]
\answerExplain{$0.45$}{Divide $9 \div 20 = 0.45$. The denominator $20 = 2^2 \times 5$ has only factors $2$ and $5$.}

\prob $-\dfrac{11}{25}$ \answerBlank[2cm]
\answerExplain{$-0.44$}{Divide $11 \div 25 = 0.44$. The denominator $25 = 5^2$. Add the negative sign: $-0.44$.}
\end{multicols}
\end{practiceBox}

% ── Practice B: Repeating Decimals ───────────────────────────────────────────
\begin{practiceBox}[funTeal]{Repeating Decimals}

\practiceHeader[funTeal]{Convert each fraction to a repeating decimal. Use bar notation.}

\prob $\dfrac{4}{9}$ \answerBlank[2cm]
\answerExplain{$0.\overline{4}$}{Divide $4 \div 9 = 0.444\ldots$ The remainder keeps being $4$, so the digit $4$ repeats: $0.\overline{4}$.}

\prob $\dfrac{5}{11}$ \answerBlank[2cm]
\answerExplain{$0.\overline{45}$}{Divide $5 \div 11 = 0.454545\ldots$ The two-digit pattern $45$ repeats: $0.\overline{45}$.}

\prob $\dfrac{2}{3}$ \answerBlank[2cm]
\answerExplain{$0.\overline{6}$}{Divide $2 \div 3 = 0.666\ldots$ The digit $6$ repeats: $0.\overline{6}$.}

\prob $-\dfrac{7}{12}$ \answerBlank[2cm]
\answerExplain{$-0.58\overline{3}$}{Divide $7 \div 12 = 0.5833\ldots$ The digit $3$ repeats after $58$. With the negative sign: $-0.58\overline{3}$.}

\prob $\dfrac{1}{6}$ \answerBlank[2cm]
\answerExplain{$0.1\overline{6}$}{Divide $1 \div 6 = 0.1666\ldots$ The digit $6$ repeats after the initial $1$: $0.1\overline{6}$.}

\prob $\dfrac{5}{6}$ \answerBlank[2cm]
\answerExplain{$0.8\overline{3}$}{Divide $5 \div 6 = 0.8333\ldots$ The digit $3$ repeats after the initial $8$: $0.8\overline{3}$.}

\end{practiceBox}

% ── Practice C: Classify ─────────────────────────────────────────────────────
\begin{sortBox}{Terminating or Repeating?}

Classify each fraction as \textbf{terminating} or \textbf{repeating} without dividing. (Hint: check the prime factors of the denominator in simplest form.)

\begin{multicols}{2}
\sortCategory[funBlue]{Terminating}
\sortCategory[funOrange]{Repeating}
\end{multicols}

\medskip
Fractions to sort: $\dfrac{3}{8}$, \quad $\dfrac{5}{9}$, \quad $\dfrac{7}{20}$, \quad $\dfrac{4}{15}$, \quad $\dfrac{1}{25}$, \quad $\dfrac{11}{30}$

\prob Classify each fraction and explain. \answerBlank[6cm]
\answerExplain{Terminating: $\frac{3}{8}, \frac{7}{20}, \frac{1}{25}$. Repeating: $\frac{5}{9}, \frac{4}{15}, \frac{11}{30}$.}{$8 = 2^3$ (only $2$s $\to$ terminates). $9 = 3^2$ (has $3$ $\to$ repeats). $20 = 2^2 \times 5$ (only $2$s and $5$s $\to$ terminates). $15 = 3 \times 5$ (has $3$ $\to$ repeats). $25 = 5^2$ (only $5$s $\to$ terminates). $30 = 2 \times 3 \times 5$ (has $3$ $\to$ repeats).}

\end{sortBox}

% ── Practice D: True or False? ───────────────────────────────────────────────
\begin{practiceBox}[funPurple]{True or False?}

\trueOrFalse{Every fraction can be written as a decimal that either terminates or repeats.}
\answerExplain{True}{This is a fundamental property of rational numbers. Long division always eventually terminates or produces a repeating remainder.}

\trueOrFalse{$\frac{1}{7} = 0.\overline{142857}$}
\answerExplain{True}{Dividing $1 \div 7 = 0.142857142857\ldots$ The six-digit block $142857$ repeats.}

\trueOrFalse{$\frac{3}{10}$ is a repeating decimal.}
\answerExplain{False}{$\frac{3}{10} = 0.3$, which terminates. The denominator $10 = 2 \times 5$ has only factors $2$ and $5$.}

\end{practiceBox}

% ── Practice E: Word Problems ────────────────────────────────────────────────
\begin{practiceBox}[funBlue]{\faGlobe~Word Problems}

\wordProblem{Three friends split a $\$10$ prize equally. Express each person's share as a decimal.}{\$}
\answerExplain{$\$3.\overline{3}$ (or approximately $\$3.33$)}{$10 \div 3 = 3.333\ldots = 3.\overline{3}$. Each person gets $\$3.\overline{3}$, which is a repeating decimal.}

\wordProblem{A bolt is $\frac{5}{16}$ inch thick. Convert this to a decimal.}{inches}
\answerExplain{$0.3125$ inches}{Divide $5 \div 16 = 0.3125$. Since $16 = 2^4$, the decimal terminates.}

\end{practiceBox}

% ── Challenge ────────────────────────────────────────────────────────────────
\begin{challengeBox}

\prob Convert $\frac{17}{99}$ to a decimal. What pattern do you notice? \answerBlank[3cm]
\answerExplain{$0.\overline{17}$}{$17 \div 99 = 0.171717\ldots = 0.\overline{17}$. Pattern: dividing by $99$ makes the numerator itself become the repeating block.}

\prob Explain why $\frac{1}{6}$ is written $0.1\overline{6}$ and not $0.\overline{16}$. \answerBlank[5cm]
\answerExplain{Only the $6$ repeats, not the $1$}{In $0.1\overline{6} = 0.1666\ldots$, the digit $1$ appears only once (in the tenths place) and does not repeat. Only the $6$ in the hundredths place onward repeats. Writing $0.\overline{16}$ would mean $0.161616\ldots$, which equals $\frac{16}{99}$, not $\frac{1}{6}$.}

\end{challengeBox}

\encouragement{Outstanding work converting fractions to decimals! You can spot terminating and repeating patterns like a pro!}
